\section{Varieties}
\label{section-varieties}

\noindent
In the Stacks project we will use the following as our definition
of a variety.

\begin{definition}
\label{definition-variety}
Let $k$ be a field. A {\it variety} is a scheme $X$ over $k$
such that $X$ is integral and the structure morphism
$X \to \Spec(k)$ is separated and of finite type.
\end{definition}

\noindent
This definition has the following drawback. Suppose that
$k'/k$ is an extension of fields. Suppose that $X$
is a variety over $k$. Then the base change
$X_{k'} = X \times_{\Spec(k)} \Spec(k')$ is
not necessarily a variety over $k'$. This phenomenon (in greater
generality) will be discussed in detail in the following sections.
The product of two varieties need not be a variety
(this is really the same phenomenon). Here is an example.

\begin{example}
\label{example-product-not-a-variety}
Let $k = \mathbf{Q}$. Let $X = \Spec(\mathbf{Q}(i))$
and $Y = \Spec(\mathbf{Q}(i))$. Then the product
$X \times_{\Spec(k)} Y$ of the varieties $X$ and $Y$
is not a variety, since it is reducible. (It is isomorphic
to the disjoint union of two copies of $X$.)
\end{example}

\noindent
If the ground field is algebraically closed however, then the
product of varieties is a variety. This follows from the results
in the algebra chapter, but there we treat much more general situations.
There is also a simple direct proof of it which we present here.

\begin{lemma}
\label{lemma-product-varieties}
\begin{slogan}
Products of varieties are varieties over algebraically closed fields.
\end{slogan}
Let $k$ be an algebraically closed field.
Let $X$, $Y$ be varieties over $k$.
Then $X \times_{\Spec(k)} Y$ is a variety over $k$.
\end{lemma}

\begin{proof}
The morphism $X \times_{\Spec(k)} Y \to \Spec(k)$ is of
finite type and separated because it is the composition of the
morphisms $X \times_{\Spec(k)} Y \to Y \to \Spec(k)$
which are separated and of finite type, see
Morphisms, Lemmas \ref{morphisms-lemma-base-change-finite-type} and
\ref{morphisms-lemma-composition-finite-type}
and
Schemes, Lemma \ref{schemes-lemma-separated-permanence}.
To finish the proof it suffices to show that $X \times_{\Spec(k)} Y$
is integral.
Let $X = \bigcup_{i = 1, \ldots, n} U_i$,
$Y = \bigcup_{j = 1, \ldots, m} V_j$ be finite affine open coverings.
If we can show that each $U_i \times_{\Spec(k)} V_j$ is integral,
then we are done by
Properties, Lemmas \ref{properties-lemma-characterize-reduced},
\ref{properties-lemma-characterize-irreducible}, and
\ref{properties-lemma-characterize-integral}.
This reduces us to the affine case.

\medskip\noindent
The affine case translates into the following algebra statement: Suppose
that $A$, $B$ are integral domains and finitely generated $k$-algebras.
Then $A \otimes_k B$ is an integral domain. To get a contradiction suppose that
$$
(\sum\nolimits_{i = 1, \ldots, n} a_i \otimes b_i)
(\sum\nolimits_{j = 1, \ldots, m} c_j \otimes d_j) = 0
$$
in $A \otimes_k B$ with both factors nonzero in $A \otimes_k B$.
We may assume that $b_1, \ldots, b_n$ are $k$-linearly
independent in $B$, and that $d_1, \ldots, d_m$ are $k$-linearly independent
in $B$. Of course we may also assume that $a_1$ and $c_1$ are nonzero
in $A$. Hence $D(a_1c_1) \subset \Spec(A)$ is nonempty. By the
Hilbert Nullstellensatz
(Algebra, Theorem \ref{algebra-theorem-nullstellensatz})
we can find a maximal ideal $\mathfrak m \subset A$ contained in
$D(a_1c_1)$ and $A/\mathfrak m = k$ as $k$ is algebraically closed.
Denote $\overline{a}_i, \overline{c}_j$ the residue classes of
$a_i, c_j$ in $A/\mathfrak m = k$. The equation above becomes
$$
(\sum\nolimits_{i = 1, \ldots, n} \overline{a}_i b_i)
(\sum\nolimits_{j = 1, \ldots, m} \overline{c}_j d_j) = 0
$$
which is a contradiction with
$\mathfrak m \in D(a_1c_1)$, the linear independence of
$b_1, \ldots, b_n$ and $d_1, \ldots, d_m$, and the fact that $B$ is a domain.
\end{proof}






